\chapter{المصايد واستدعاءات النظام}
\label{chap:trap}

تعتبر \indextext{المصيدة (Trap)} الوسيلة الأساسية التي ينقل بها المعالج التحكم من مساحة المستخدم إلى مساحة النواة.
تحدث المصيدة في الحالات التالية:
\begin{enumerate}
\item صدور \indextext{استدعاء نظام (System Call)} عند تنفيذ التعليمية \texttt{ecall}.
\item وقوع \indextext{استثناء (Exception)} كالقسمة على الصفر أو خطأ الوصول للذاكرة.
\item وصول \indextext{مقاطعة أجهزة (Device Interrupt)} من أحد الملحقات الخارجية.
\end{enumerate}

\section{معالجة المصايد في RISC-V}

عند وقوع مصيدة بداخل وضع المستخدم، ينجز عتاد المعالج في RISC-V الخطوات التلقائية التالية:
\begin{enumerate}
\item يحفظ عنوان التعليمية الحالية في السجل \texttt{sepc}.
\item يحفظ سبب وقوع المصيدة في السجل \texttt{scause}.
\item يحول وضع المعالج من وضع المستخدم U-mode إلى وضع المشرف S-mode.
\item ينسخ عنوان معالج المصايد المخزن بداخل السجل \texttt{stvec} إلى عداد البرنامج \texttt{pc} لينتقل التنفيذ إلى النواة.
\end{enumerate}

يتولى الملف المجمع \texttt{trampoline.S} حفظ كامل سجلات العملية بداخل صفحة المجمّع المخصصة (Trapframe)،
ثم يبدل السجل \texttt{satp} ليشير إلى جدول صفحات النواة، وينتقل لتنفيذ الدالة \texttt{usertrap()} بداخل الملف \texttt{trap.c}.

تتحقق الدالة \texttt{usertrap()} من سبب المصيدة عبر السجل \texttt{scause}؛
فإذا كان السبب استدعاء نظام، تستدعي الدالة \texttt{syscall()} التي تقرأ رقم الاستدعاء من السجل \texttt{a7} وتنفذ الدالة المطلوبة،
ثم تعيد نتيجة التنفيذ بداخل السجل \texttt{a0}.
